\chapter{Safety}\label{ch:safety}

Read this chapter before you connect power to the machine. The servos of a
plotter are strong and fast. They can cause injury when they move.

\section{Status of the Software}

\begin{warning}
ClearHPGL is experimental. It is not certified for production or for
applications where safety is critical. Do not use it as the only safety
function of a machine. Install a physical emergency stop that removes power
from the motors.
\end{warning}

The emergency stop input of the board (\conn{DI-6}) stops the job and
disables the motors through the firmware. Firmware can fail. A hardware
emergency stop in the motor power circuit does not depend on the firmware.
You must have both.

\section{Motion That You Do Not Expect}

\begin{warning}
The motors enable automatically when a job arrives (setting
\key{auto_enable}~1). A job from a computer on the network can start the
machine at any time. Keep your hands and tools out of the work area when the
board has power and a computer is connected.
\end{warning}

\begin{warning}
Without limit switches, the machine has no protection at the ends of its
travel. It goes to the coordinates that the job gives, inside the plot area
that you set. If the plot area is larger than the mechanics, the carriage can
hit the end stops. Set \key{width_mm} and \key{height_mm} to the travel that
is really available. Limit switches (chapter~\ref{ch:limits}) stop the axis
with a ramp: they are not an emergency stop.
\end{warning}

\begin{itemize}
  \item Before you plot, look at the preview in the desktop app. Make sure
        that the drawing is in the plot area.
  \item Before you change the steps per revolution, the travel per
        revolution, or the direction of an axis, do a test move at low
        speed. Chapter~\ref{ch:setup} gives the procedure.
  \item Keep the work area clear of persons who do not operate the machine.
\end{itemize}

\section{Faults and the Emergency Stop}\label{sec:fault-behavior}

The board stops the machine when one of these conditions occurs:

\begin{itemize}
  \item The emergency stop input \conn{DI-6} opens (with the default setting
        \key{estop_di6}~1).
  \item A ClearPath servo reports an alert or a fault.
  \item The ready signal (HLFB) of a servo is not present when the motors
        enable (with \key{hlfb_check}~1).
\end{itemize}

When a fault occurs, the board disables the motors, drops the remainder of
the job, and latches the fault. The servos stop without control of
deceleration. The pen can stay on the material. To continue, remove the
cause, then clear the fault with the \cmd{clear} command or with
\menu{Clear fault} in the desktop app.

\begin{caution}
When the motors disable, the axes are free. On a vertical machine or a
machine on a slope, gravity can move an axis. Make sure that the axes cannot
fall or move when the motors have no power.
\end{caution}

\begin{note}
\notelabel If you do not connect an emergency stop to \conn{DI-6}, set
\key{estop_di6} to 0. If you do not do this, the board stays in the fault
state and does not plot. A hardware emergency stop in the motor power circuit
is still necessary.
\end{note}

\section{Lasers}

\begin{warning}
Do not connect a laser to the pen output. The pen output turns on and off
with the timing of a pen: it has settle delays and no control of power. The
firmware has no interlocks for an enclosure, a lid, or air assist. A laser on
this output can cause fire, burns, and damage to the eyes.
\end{warning}

Laser control is a plan for a future version. It will need control of the
power by the speed of the path, and interlocks that can stop the beam.

\section{Electrical Safety}

\begin{warning}
Disconnect all power before you connect or disconnect the motors, the pen
lift, or the emergency stop. The ClearPath servos and their power supplies
have dangerous voltages. Obey the instructions of Teknic for the hardware.
\end{warning}

\begin{caution}
The board writes its settings to nonvolatile memory. It can do this only
when the supply voltage is above the undervoltage limit of the ClearCore. Do
not change settings while the ClearCore gets its power only from the USB
cable.
\end{caution}
